\documentclass[10pt,aspectratio=169]{beamer}
\input{../../_en_preambulo_beamer}
% Comandos y macros estadísticas compartidas para las presentaciones Beamer en inglés (Metropolis)

% Conjuntos numéricos básicos
\newcommand{\R}{\mathbb{R}}
\newcommand{\N}{\mathbb{N}}
\newcommand{\Q}{\mathbb{Q}}
\newcommand{\Z}{\mathbb{Z}}
\newcommand{\set}[1]{\left\{ #1 \right\}}
\newcommand{\abs}[1]{\left|#1\right|}
\newcommand{\norm}[1]{\left\| #1 \right\|}

% Comandos estadísticos y probabilísticos del curso
\newcommand{\Var}{\operatorname{Var}}
\newcommand{\cov}{\operatorname{Cov}}
\newcommand{\corr}{\rho}
\newcommand{\comb}[2]{\begin{pmatrix} #1 \\ #2 \end{pmatrix}}
\newcommand{\s}{\sigma}
\newcommand{\particion}{\mathfrak{P}}
\newcommand{\card}[1]{\#\left( #1 \right)}
\newcommand{\seq}[3]{\set{#1}_{#2}^{#3}}
\newcommand{\E}{\mathbb{E}}
\newcommand{\Prob}{\mathbb{P}}

% Entornos pedagógicos y cajas en inglés académico natural (soportando tanto nombres en inglés como en español)
\newenvironment{discussion}{%
  \begin{alertblock}{Discussion Question}%
}{%
  \end{alertblock}%
}
\newenvironment{discusion}{%
  \begin{alertblock}{Discussion Question}%
}{%
  \end{alertblock}%
}

\newenvironment{reflection}{%
  \begin{alertblock}{Classroom Reflection}%
}{%
  \end{alertblock}%
}
\newenvironment{reflexion}{%
  \begin{alertblock}{Classroom Reflection}%
}{%
  \end{alertblock}%
}

\newenvironment{paradox}{%
  \begin{alertblock}{Exploratory Paradox}%
}{%
  \end{alertblock}%
}
\newenvironment{paradoja}{%
  \begin{alertblock}{Exploratory Paradox}%
}{%
  \end{alertblock}%
}

\newenvironment{motivation}{%
  \begin{exampleblock}{Motivation: Data Science in Practice}%
}{%
  \end{exampleblock}%
}
\newenvironment{motivacion}{%
  \begin{exampleblock}{Motivation: Data Science in Practice}%
}{%
  \end{exampleblock}%
}

\newenvironment{quickchallenge}{%
  \begin{block}{Quick Team Challenge}%
}{%
  \end{block}%
}
\newenvironment{retoclase}{%
  \begin{block}{Quick Team Challenge}%
}{%
  \end{block}%
}


\title[Z and t Statistics]{Section 5.9: Z and t Statistics}
\subtitle{Standardization with known or unknown standard deviation}
\author[J. Castillo Colmenares]{Juliho Castillo Colmenares}
\institute{Tecnológico de Monterrey}
\date{}

\begin{document}
\begin{frame}[plain]\vspace{-0.8cm}\titlepage\end{frame}

\begin{frame}{Roadmap --- Chapter 5}
  \footnotesize
  \begin{itemize}\itemsep=0.08cm
    \item \textbf{05.08} Fundamental concepts
    \item \textbf{05.09} Z and t statistics \emph{(today)}
    \item \textbf{06--07} Confidence intervals and hypothesis tests
  \end{itemize}
  \begin{block}{Learning objective}
    Construct $Z$ and $t$, distinguish known from unknown $\sigma$, and read
    quantiles of Student's $t$ distribution.
  \end{block}
\end{frame}

\begin{frame}{Inferential setup}
  \footnotesize
  The reading starts with a null-hypothesis value $A_0$ and a random sample of
  size $n$.
  \begin{block}{Observed quantity}
    The sample mean or measure is denoted $A$ and compared with $A_0$.
  \end{block}
\end{frame}

\begin{frame}{Z statistic}
  \footnotesize
  \[
    Z=\frac{A-A_0}{\sigma/\sqrt n}.
  \]
  Here $\sigma$ is the known population standard deviation.
  \begin{alertblock}{Purpose}
    Convert a normal variable into a standard normal variable for probabilities
    and tabulated quantiles.
  \end{alertblock}
\end{frame}

\begin{frame}{Z case: known $\sigma$}
  \footnotesize
  If the standard deviation is known from previous experience, the Z statistic
  uses $\sigma/\sqrt n$ directly.
  \begin{exampleblock}{Reading context}
    Pizza delivery times illustrate a mean whose population variability is
    documented.
  \end{exampleblock}
\end{frame}

\begin{frame}{t case: unknown $\sigma$}
  \footnotesize
  When no reliable record of $\sigma$ exists or the sample is small, replace it
  with the sample standard deviation $S$:
  \[
    t=\frac{A-A_0}{S/\sqrt n}.
  \]
  \begin{alertblock}{Difference}
    The extra uncertainty is represented by Student's $t$ distribution.
  \end{alertblock}
\end{frame}

\begin{frame}{Student's distribution}
  \footnotesize
  Student's $t$ distribution was described by William Sealy Gosset under the
  pseudonym Student. It is symmetric about zero and its tails depend on the
  degrees of freedom.
  \begin{block}{Use}
    It models standardization when the population deviation is estimated from
    the sample itself.
  \end{block}
\end{frame}

\begin{frame}{Degrees of freedom and variance}
  \footnotesize
  The parameter `df` is denoted $\nu$. For $X\sim t_\nu$,
  \[
    \mu_X=0,\qquad \sigma_X^2=\frac{\nu}{\nu-2}
    \quad(\nu>2).
  \]
  As $\nu$ grows, $t_\nu$ approaches the standard normal distribution.
\end{frame}

\begin{frame}{Sample variance and t statistic}
  \footnotesize
  The reading presents
  \[
    S^2=\sum\frac{(A_i-A_0)^2}{n-1},
    \qquad t=\frac{A-A_0}{S/\sqrt n}.
  \]
  \begin{alertblock}{Reading}
    The divisor $n-1$ and the degrees of freedom connect estimation of $S$ with
    the shape of Student's distribution.
  \end{alertblock}
\end{frame}

\begin{frame}{Computational bridge 1/4: choose Z or t}
  \footnotesize
  \begin{tabular}{lll}
    \hline
    Available information & Statistic & Reference distribution \\
    \hline
    $\sigma$ known & $Z=(A-A_0)/(\sigma/\sqrt n)$ & standard normal \\
    $\sigma$ unknown & $t=(A-A_0)/(S/\sqrt n)$ & $t_\nu$ \\
    \hline
  \end{tabular}
\end{frame}

\begin{frame}{Computational bridge 2/4: degrees of freedom}
  \footnotesize
  Degrees of freedom determine the tails of $t_\nu$. With less information
  about the standard deviation, the tails are heavier.
  \begin{block}{Limit}
    $t_\nu\to N(0,1)$ as $\nu$ grows.
  \end{block}
\end{frame}

\begin{frame}{Computational bridge 3/4: areas}
  \footnotesize
  For an observed value $t$, the right-tail probability is
  \[
    P(T_\nu>t)=1-F_{t_\nu}(t).
  \]
  Symmetry gives left-tail quantiles by changing the sign.
\end{frame}

\begin{frame}{Computational bridge 4/4: quantiles}
  \footnotesize
  The percent point function satisfies
  \[
    F_{t_\nu}(t_p)=p.
  \]
  A right-tail area of $0.05$ uses $p=0.95$ and symmetry for the negative limit.
\end{frame}

\begin{frame}{Solved example 1/4 --- $\nu=9$ quantiles}
  \footnotesize
  Consider a $t$ variable with $\nu=9$ degrees of freedom. Find values with a
  right-tail area of $0.05$ and a central unshaded area of $0.90$.
  \begin{block}{Statement}
    Find the two symmetric limits.
  \end{block}
\end{frame}

\begin{frame}{Solved example 1/4 --- solution}
  \footnotesize
  The left area of the positive limit is $0.95$:
  \[
    t_{0.95,9}\approx1.833,
    \qquad t_{0.05,9}\approx-1.833.
  \]
  \begin{alertblock}{Check}
    $P(-1.833<T<1.833)=0.95-0.05=0.90$.
  \end{alertblock}
\end{frame}

\begin{frame}{Reading application 2/4 --- Z test}
  \footnotesize
  A delivery process has a known population standard deviation. The observed
  mean is $A$ and the reference value is $A_0$.
  \begin{block}{Statement}
    Standardize the difference using the section's Z formula.
  \end{block}
\end{frame}

\begin{frame}{Reading application 2/4 --- Z interpretation}
  \footnotesize
  \[
    Z=\frac{A-A_0}{\sigma/\sqrt n}.
  \]
  A positive value means $A$ lies above $A_0$; its magnitude is read on a
  standard normal distribution.
\end{frame}

\begin{frame}{Reading application 3/4 --- t test}
  \footnotesize
  If $\sigma$ is unavailable, compute $S$ from the sample and use
  \[
    t=\frac{A-A_0}{S/\sqrt n}.
  \]
  \begin{block}{Statement}
    The uncertainty about $\sigma$ must be represented by $t_\nu$.
  \end{block}
\end{frame}

\begin{frame}{Reading application 3/4 --- t interpretation}
  \footnotesize
  Compare the observed value with $t_\nu$ quantiles rather than normal
  quantiles, especially for a small $n$.
  \begin{alertblock}{Result}
    Heavier tails protect against underestimating variability.
  \end{alertblock}
\end{frame}

\begin{frame}{Reading application 4/4 --- sample variance}
  \footnotesize
  The section defines $S^2$ with divisor $n-1$ and then forms the t statistic
  with $S$.
  \begin{block}{Statement}
    Identify which quantities change as the sample size grows.
  \end{block}
\end{frame}

\begin{frame}{Reading application 4/4 --- connection}
  \footnotesize
  As $n$ grows, degrees of freedom increase, $S$ estimates $\sigma$ more closely,
  and $t_\nu$ approaches $N(0,1)$.
  \begin{exampleblock}{Integrated reading}
    Formula, sample variance, and reference distribution form one standardizing
    procedure.
  \end{exampleblock}
\end{frame}

\begin{frame}{Synthesis and transition}
  \footnotesize
  \begin{itemize}\itemsep=0.12cm
    \item Z uses known $\sigma$ and a normal reference.
    \item t uses $S$, degrees of freedom, and Student tails.
    \item Symmetry and the percent point function provide area limits.
  \end{itemize}
  \begin{alertblock}{Next step}
    These tools are applied in the confidence-interval and hypothesis-testing
    sections.
  \end{alertblock}
\end{frame}
\end{document}
